An observer in Deh-Namak (you will be taking the observational part of the
exam in this region tonight) has observed Venus for seven months, started
from September 2008 and continued until March 2009. During the observation, a
research CCD camera and an image processing software were used to take high
resolution images and to extract high precision data. Table 2.1 shows the
collected data during the observation.

Table 2.1 description:
\begin{center}
\begin{tabular}{lp{9cm}}
Column 1 & Date of observation. \\
Column 2 & Earth--Sun distance in astronomical unit (AU) for observation
  date and time. This value is taken from high precision tables. \\
Column 3 & Phase of Venus, Percent of Venus disk illuminated by the Sun as
  observed from the Earth. \\
Column 4 & Elongation of Venus, the angular distance between center of the
  Sun and center of Venus in degrees as observed from the Earth. \\
\end{tabular}
\end{center}

\begin{center}
Table 2.1

\begin{tabular}{c|c|c|c}
Column 1 & Column 2 & Column 3 & Column 4 \\
\hline
Date & Earth--Sun Distance (AU) & Phase (\%)
  & Elongation (SEV) ($^\circ$) \\
\hline
20/9/2008   & 1.0043 & 88.4 & 27.56 \\
10/10/2008  & 0.9986 & 84.0 & 32.29 \\
20/10/2008  & 0.9957 & 81.6 & 34.53 \\
30/10/2008  & 0.9931 & 79.0 & 36.69 \\
9/11/2008   & 0.9905 & 76.3 & 38.71 \\
19/11/2008  & 0.9883 & 73.4 & 40.62 \\
29/11/2008  & 0.9864 & 70.2 & 42.38 \\
19/12/2008  & 0.9839 & 63.1 & 45.29 \\
29/12/2008  & 0.9834 & 59.0 & 46.32 \\
18/1/2009   & 0.9838 & 49.5 & 47.09 \\
7/2/2009    & 0.9863 & 37.2 & 44.79 \\
17/2/2009   & 0.9881 & 29.6 & 41.59 \\
27/2/2009   & 0.9904 & 20.9 & 36.16 \\
19/3/2009   & 0.9956 & 3.8  & 16.08 \\
\end{tabular}
\end{center}

\begin{description}
\item[a)] Using given data in table 2.1, calculate the Sun-Venus-Earth angle
  ($\angle SVE$). This is angular separation between the Sun and the Earth as
  seen from Venus. Write $\angle SVE$ angle in column 2 of Table 2.2 in your
  answer sheet for the all observing dates.

  Hint: Remember that the line between light and shadow, in the phases, is an
  ellipse arc.

\ioaafig{2009_DA_D2-f1.pdf}
\begin{center}
Figure (part a)
\end{center}

\item[b)] Calculate Sun--Venus distance in AU and write it down in column 3
  of table 2.2 for all observation dates.
\item[c)] Plot Sun--Venus distance versus observing date.
\item[d)] Find perihelion ($r_{v,min}$) and aphelion ($r_{v,max}$) distances
  of Venus from the Sun.
\item[e)] Calculate semi-major axis ($a$) of the Venus orbit.
\item[f)] Calculate eccentricity ($e$) of Venus orbit.
\end{description}
